\subsection{Formal Definitions}
\label{sec:formal_definitions}

\begin{definition}[SemanticFact]
\label{def:semanticfact}
A \emph{semantic fact} is a tuple
\[
f = (i, s, r, \vec{o}, c, k, m)
\]
where:
\begin{itemize}
    \item $i \in \mathcal{I}$ is a globally unique identifier.
    \item $s \in \mathcal{E}$ is the subject entity.
    \item $r \in \mathcal{R}$ is the relation type.
    \item $\vec{o} \in \mathcal{V}^*$ is an ordered sequence of values (arity $\operatorname{arity}(f) = |\vec{o}|$).
    \item $c \in \mathcal{C}$ is the context of validity.
    \item $k \in [0,1]$ is the confidence.
    \item $m \in \mathbf{Set}$ is metadata.
\end{itemize}
Let $\mathcal{F}$ denote the set of all semantic facts.
\end{definition}

$\mathcal{E} \subseteq \mathcal{I}$ and $\mathcal{R} \subseteq \mathcal{I}$
denote the sets of entities and relations respectively, each carrying a
type and (for relations) an arity that a fact's object tuple must match;
Appendix~\ref{sec:appendix:definitions} gives the full structure, which is
not needed to follow the rest of the paper.

\begin{definition}[Context]
\label{def:context}
A \emph{context} $c \in \mathcal{C}$ is a node in a finite poset $(\mathcal{C}, \leq)$
identified by a dot-separated path. The partial order is prefix inclusion:
$c_1 \leq c_2$ iff $c_2$ is a prefix of $c_1$ (i.e., $c_1$ is more specific).
The root context $\top$ (``world'') is the unique maximal element:
$c \leq \top$ for all $c \in \mathcal{C}$.
Joins always exist: $c_1 \vee c_2$ (the least upper bound) is the longest
common prefix, i.e.\ the most specific common ancestor. Meets are partial:
if $c_1 \leq c_2$ then $c_1 \wedge c_2 = c_1$, but two incomparable
contexts have no common lower bound at all, since no context can carry two
distinct sibling paths as prefixes. The context hierarchy is therefore a
tree-shaped join-semilattice, not a lattice
(Appendix~\ref{sec:appendix:theorems} states the resulting open-set
properties).
\end{definition}

\begin{definition}[Finite Topological Space]
\label{def:topology}
A \emph{finite topological space} is a pair $(X, \mathcal{T})$ where $X$ is a finite set
and $\mathcal{T} \subseteq \mathcal{P}(X)$ satisfies:
\begin{enumerate}
    \item $\emptyset, X \in \mathcal{T}$.
    \item $U, V \in \mathcal{T} \implies U \cap V \in \mathcal{T}$ (finite intersection).
    \item $\{U_i\}_{i \in I} \subseteq \mathcal{T} \implies \bigcup_i U_i \in \mathcal{T}$ (arbitrary union).
\end{enumerate}
In SFDB, $X = \mathcal{I}$ (fact identifiers) and $\mathcal{T}$ is generated
by the Alexandrov topology on $(\mathcal{C}, \leq)$.
\end{definition}

\paragraph{Remark.}
Section~\ref{sec:sheaf-remark}, immediately below, explains why the sheaf
condition holds vacuously on this topology and is not the paper's
technical contribution; we do not repeat that argument here.

An \emph{open set} $U \in \mathcal{T}$ is a subset of $X$ declared open; in
the Alexandrov topology induced by $(\mathcal{C}, \leq)$, every down-set
$\downarrow c = \{d \in \mathcal{C} \mid d \leq c\}$ is open, and each
context $c$ has an associated open set
$U_c = \{f \in \mathcal{I} \mid \operatorname{context}(f) \leq c\}$
(Appendix~\ref{sec:appendix:definitions} gives this formally, together
with its intersection and top-element properties).

\begin{definition}[Neighbourhood]
\label{def:neighbourhood}
A \emph{neighbourhood} of $x \in X$ is any $U \in \mathcal{T}$ with $x \in U$.
The neighbourhood filter is $\mathcal{N}(x) = \{U \in \mathcal{T} \mid x \in U\}$.
In the Alexandrov topology, the minimal neighbourhood exists uniquely:
$\mathcal{B}(x) = \bigcap_{U \in \mathcal{N}(x)} U = \downarrow \operatorname{context}(x)$,
using the same context-assignment function $\operatorname{context}: \mathcal{I} \to \mathcal{C}$
that defines $U_c$ above.
\end{definition}

\begin{definition}[Restriction Map]
\label{def:restriction}
For $c \leq d$ (so $U_c \subseteq U_d$), the \emph{restriction map}
$\rho_{d,c}: F(d) \to F(c)$ is the identity on facts, restricted to the
subset of $F(d)$ that already lies in $U_c$:
\[
\rho_{d,c}(f) =
\begin{cases}
f & \text{if } \operatorname{context}(f) \leq c, \\
\text{undefined} & \text{otherwise.}
\end{cases}
\]
A fact is not rewritten by restriction --- no field of $f$ changes --- it
is simply re-indexed as an element of the smaller set $F(c)$ when its own
context already places it there. This is the presheaf of \emph{visibility
from an open set}: restricting a broader view to a narrower one either
finds the fact was already visible from that narrower vantage point, or
it was not, in which case restriction is undefined rather than
approximate. (An earlier version of this definition described
restriction as dropping object slots not meaningful in the target
context; the implementation does not perform any such transformation,
and that description is retracted here rather than left standing
alongside code that does something else.)
Functoriality: $\rho_{c,c} = \operatorname{id}_{F(c)}$, and for
$c \leq d \leq e$, $\rho_{e,c} = \rho_{d,c} \circ \rho_{e,d}$ holds
because both sides equal the identity restricted to
$\{f \in F(e) : \operatorname{context}(f) \leq c\}$.
\end{definition}

\begin{definition}[Stalk]
\label{def:stalk}
The \emph{stalk} of presheaf $F$ at point $x \in X$ is
$F_x = \varinjlim_{U \in \mathcal{N}(x)} F(U)$,
the set of germs $[s]_U$ where $[s]_U \sim [t]_V$ iff $\exists W \subseteq U \cap V$
with $x \in W$ and $s|_W = t|_W$.
\end{definition}

On an Alexandrov space this direct limit is not an abstract object: every
point has a minimal open neighbourhood $\mathcal{B}(x)$
(Definition~\ref{def:neighbourhood}), the limit collapses to
$F_x \cong F(\mathcal{B}(x))$, and no germ equivalence needs to be
computed. This is what licenses the implementation's \emph{stalk index}:
a hash table mapping each fact identifier $x$ to the section over its
minimal neighbourhood is literally a table of stalks, materialised in
advance, and the paper's ``stalk lookup'' is a lookup in that table
rather than a direct-limit computation at query time.

\begin{definition}[Local Section]
\label{def:localsection}
A \emph{local section} is $s \in F(U)$ for a proper open subset $U \subsetneq X$.
In SFDB, a local section is a fact whose context $c < \top$.
\end{definition}

\begin{definition}[Global Section]
\label{def:globalsection}
A \emph{global section} is $s \in F(X)$ over the whole space.
In SFDB, this is a fact valid in the root context $\top$.
Computed by collecting local sections, checking compatibility on overlaps,
and gluing upward until $\top$ is reached.
\end{definition}

A \emph{knowledge state} $\Sigma = (X, F)$ pairs a finite topological space
of fact identifiers with a presheaf $F$ assigning facts to open sets, with
insert/delete/update as set operations on $\Sigma$; a \emph{query}
$Q = (\tau, s, r, \vec{o}, c, k_{\min}, n_{\max})$ denotes the set of facts
$\llbracket Q \rrbracket_\Sigma$ matching its subject/relation/object/context/confidence
pattern. Both are given fully in Appendix~\ref{sec:appendix:definitions}.
Section~\ref{sec:correctness} defines the notion of \emph{result
equivalence} actually used by the implementation and the equivalence
proof (Definition~\ref{def:semanticequiv}); it is the only such
definition in this paper.
